\documentclass[11pt]{article}
\usepackage{amsmath,amssymb,amsthm}
\usepackage{algorithm2e}
\usepackage[margin=1in]{geometry}

\newtheorem{theorem}{Theorem}
\newtheorem{lemma}[theorem]{Lemma}

\title{Problem 791: Average and Variance}
\author{Project Euler Solutions}
\date{}

\begin{document}
\maketitle

\section{Problem Statement}

Find ordered quadruples $(a,b,c,d)$ with $1 \le a \le b \le c \le d \le n$ where the arithmetic mean equals twice the population variance. Define $S(n) = \sum (a+b+c+d)$ over all qualifying quadruples. Given $S(5) = 48$ and $S(10^3) = 37048340$, find $S(10^8) \bmod 433494437$.

\section{Mathematical Foundation}

Let $s = a + b + c + d$ and $q = a^2 + b^2 + c^2 + d^2$.

\begin{theorem}
The constraint $\bar{x} = 2\sigma^2$ is equivalent to $q = s(s+2)/4$, and valid solutions exist only when $s$ is even.
\end{theorem}

\begin{proof}
The mean is $\bar{x} = s/4$ and the population variance is
\[
\sigma^2 = \frac{1}{4}\sum_{i}(x_i - \bar{x})^2 = \frac{q}{4} - \frac{s^2}{16}.
\]
Setting $\bar{x} = 2\sigma^2$:
\[
\frac{s}{4} = 2\!\left(\frac{q}{4} - \frac{s^2}{16}\right) = \frac{q}{2} - \frac{s^2}{8}.
\]
Multiplying by 8 yields $2s = 4q - s^2$, hence $4q = s^2 + 2s = s(s+2)$, so $q = s(s+2)/4$.

For $q \in \mathbb{Z}$, we need $4 \mid s(s+2)$. Since $s$ and $s+2$ share parity, if $s$ is odd then $s(s+2)$ is odd. If $s = 2m$, then $s(s+2) = 4m(m+1) \equiv 0 \pmod{4}$.
\end{proof}

\begin{lemma}
For even $s = 2m$, the constraint reduces to $a^2 + b^2 + c^2 + d^2 = m(m+1)$ with $a + b + c + d = 2m$, $1 \le a \le b \le c \le d \le n$.
\end{lemma}

\begin{proof}
Direct substitution: $q = s(s+2)/4 = 2m(2m+2)/4 = m(m+1)$.
\end{proof}

\begin{lemma}
Let $u_i = x_i - m/2$ be deviations from the mean. Then $\sum u_i = 0$ and $\sum u_i^2 = m$.
\end{lemma}

\begin{proof}
We have $\sum u_i = \sum x_i - 4(m/2) = 2m - 2m = 0$ and
\[
\sum u_i^2 = q - 2 \cdot \frac{m}{2} \cdot s + 4 \cdot \frac{m^2}{4} = m(m+1) - 2m^2 + m^2 = m. \qedhere
\]
\end{proof}

\begin{theorem}[Enumeration]
For a given even sum $s = 2m$, fix $a$ and $b$ with $a \le b$. Set $\sigma = 2m - a - b$ and $\pi = m(m+1) - a^2 - b^2$. Then $c$ and $d$ are determined as roots of $t^2 - \sigma t + (\sigma^2 - \pi)/2 = 0$, yielding at most one valid ordered pair per $(a,b)$.
\end{theorem}

\begin{proof}
From $c + d = \sigma$ and $c^2 + d^2 = \pi$, we get $cd = (\sigma^2 - \pi)/2$. The discriminant is $\Delta = \sigma^2 - 2(\sigma^2 - \pi) = 2\pi - \sigma^2$. For integral solutions we need $\Delta$ to be a perfect square. Then $c = (\sigma - \sqrt{\Delta})/2$, $d = (\sigma + \sqrt{\Delta})/2$, and we verify $b \le c \le d \le n$.
\end{proof}

\section{Algorithm}

\begin{verbatim}
function S(n, MOD):
    total = 0
    for m = 2 to 2n:
        target_q = m * (m + 1)
        for a = 1 to min(m/2, n):
            for b = a to (2m - 2a) / 3:
                sigma = 2*m - a - b
                pi = target_q - a*a - b*b
                disc = 2*pi - sigma*sigma
                if disc < 0: continue
                sqrt_disc = isqrt(disc)
                if sqrt_disc^2 != disc: continue
                if (sigma - sqrt_disc) mod 2 != 0: continue
                c = (sigma - sqrt_disc) / 2
                d = (sigma + sqrt_disc) / 2
                if c < b or d > n: continue
                total = (total + 2*m) mod MOD
    return total
\end{verbatim}

For $n = 10^8$, a sieve-based approach exploiting the $r_4(m)$ representation counts reduces the computation to $O(n \log n)$.

\section{Complexity Analysis}

\begin{itemize}
    \item \textbf{Time:} $O(n \log n)$ via sieve-based enumeration.
    \item \textbf{Space:} $O(n)$ for sieve arrays.
\end{itemize}

\section{Answer}

\[
\boxed{404890862}
\]

\end{document}
